\subsection{Eventy — implementácia}

\subsubsection{Tvorba eventu}

Obrazovka tvorby eventu (\texttt{EventCreationScreen.kt}) sa otvorí po potvrdení
polohy na mape. Formulár obsahuje nasledujúce polia:

\begin{itemize}
    \item \textbf{Názov a popis} — textové polia, obe sú povinné.
    \item \textbf{Miesto} — predvyplnené adresou získanou z Geocoderu / Places API.
    \item \textbf{Cena a počet účastníkov} — hodnota 0 znamená zadarmo, resp. neobmedzený počet.
    \item \textbf{Dátum a čas} — výber cez \texttt{DatePickerDialog} a \texttt{TimePickerDialog}.
    \item \textbf{Trvanie} — zadávané v hodinách a minútach; hodnota \texttt{dateTo}
          sa vypočíta automaticky:
\begin{lstlisting}[language=Kotlin]
val computedDateTo = dateFrom?.let {
    Date(it.time + totalMinutes * 60 * 1000L)
}
\end{lstlisting}
    \item \textbf{Kategória a podkategória} — výber z rozbaľovacieho zoznamu
          definovaného v \texttt{CATEGORY\_MAP}.
    \item \textbf{Viditeľnosť} — verejné alebo súkromné; súkromný event vyžaduje heslo.
    \item \textbf{Typ check-in} — jedna zo štyroch možností:
          \texttt{none}, \texttt{qr}, \texttt{manual}, \texttt{qr\_and\_manual}.
    \item \textbf{Deadline odhlásenia} — počet hodín pred eventom, po uplynutí ktorých
          sa odhlásenie penalizuje.
\end{itemize}

\paragraph{Upload obrázkov cez Cloudinary:}
Obrázky sa vyberajú z galérie zariadenia pomocou
\texttt{ActivityResultContracts.GetMultipleContents()}. Každý obrázok sa uploaduje
ako \texttt{multipart/form-data} priamo na Cloudinary API bez potreby vlastného
backendu:

\begin{lstlisting}[language=Kotlin, caption=Upload obrázka na Cloudinary]
private suspend fun uploadToCloudinary(
    context: Context,
    uri: Uri,
    eventId: String
): String {
    val url = URL("https://api.cloudinary.com/v1_1/$cloudName/image/upload")
    val conn = (url.openConnection() as HttpsURLConnection).apply {
        requestMethod = "POST"
        setRequestProperty(
            "Content-Type",
            "multipart/form-data; boundary=$boundary"
        )
    }
    // zapíše upload_preset + bytes obrázka
    // vráti secure_url z JSON odpovede
}
\end{lstlisting}

Prístupové kľúče (\texttt{CLOUDINARY\_CLOUD\_NAME}, \texttt{CLOUDINARY\_UPLOAD\_PRESET})
sú uložené v súbore \texttt{local.properties} a sprístupnené cez \texttt{BuildConfig}.

\paragraph{Uloženie eventu:}
Po vyplnení formulára sa event uloží volaním \texttt{onSave()}:

\begin{lstlisting}[language=Kotlin, caption=Uloženie eventu]
onSave(event.copy(
    title               = title.trim(),
    dateFrom            = dateFrom,
    dateTo              = computedDateTo,
    visibility          = if (isPrivate) "private" else "public",
    imageUrls           = existingUrls + uploadedUrls,
    checkInType         = checkInType,
    cancelDeadlineHours = cancelDeadlineInput.toIntOrNull() ?: 0
))
\end{lstlisting}

Metóda \texttt{AppViewModel.saveEventToMap()} uloží event súčasne do in-memory
mapy \texttt{\_events} aj do databázy Firestore:

\begin{lstlisting}[language=Kotlin, caption=Uloženie eventu do Firestore]
fun saveEventToMap(event: Event) {
    _events[event.id] = event
    viewModelScope.launch {
        db.collection("events")
          .document(event.id)
          .set(event.toMap())
          .await()
    }
}
\end{lstlisting}

% -------------------------------------------------------
\subsubsection{Detail eventu}

Obrazovka detailu eventu (\texttt{EventDetailScreen.kt}) je rozdelená do troch tabov.

\paragraph{Tab 0 — Detail:}
Zobrazuje všetky informácie o evente. Kľúčové stavy používateľa sú určené
nasledovne:

\begin{lstlisting}[language=Kotlin, caption=Stavové premenné v detaile eventu]
val isCreator    = event.creatorId == currentUserId
val isAttending  = currentUserId in event.attendees
val isOnWaitlist = currentUserId in event.waitlist
val isFull       = event.participants > 0 &&
                   event.attendees.size >= event.participants
val isBanned     = currentUserId in event.bannedUsers
\end{lstlisting}

Tlačidlá akcie sa dynamicky menia podľa stavu — tvorca eventu vidí možnosti
úpravy a mazania, prihlásený účastník vidí odhlásenie, ostatní používatelia
vidia možnosť prihlásenia alebo zaradenia do poradovníka.

\paragraph{Tab 1 — Účastníci:}
Zobrazuje zoznam účastníkov s ich avatarmi. Tvorca eventu má navyše možnosť:
\begin{itemize}
    \item manuálne označiť príchod účastníka (ak \texttt{checkInType == "manual"}),
    \item odstrániť účastníka z eventu — odstránený účastník je zároveň zaradený
          do zoznamu \texttt{bannedUsers} a prvý používateľ z poradovníka je
          automaticky povýšený.
\end{itemize}

\paragraph{Tab 2 — Chat:}
Real-time chat implementovaný pomocou Firestore \texttt{snapshotListener}.
Prístup do chatu majú výlučne prihlásení účastníci a tvorca eventu.

% -------------------------------------------------------
\subsubsection{Správa účasti}

Správa účasti je implementovaná v \texttt{AppViewModel} prostredníctvom
troch metód.

\paragraph{joinEvent — prihlásenie na event:}
Ak je event plný, používateľ je zaradený do poradovníka; v opačnom prípade
priamo medzi účastníkov:

\begin{lstlisting}[language=Kotlin, caption=Prihlásenie na event]
fun joinEvent(eventId: String, userId: String) {
    val isFull = event.participants > 0 &&
                 event.attendees.size >= event.participants
    val updated = if (isFull)
        event.copy(waitlist  = event.waitlist  + userId)
    else
        event.copy(attendees = event.attendees + userId)
    // uloží do Firestore
}
\end{lstlisting}

\paragraph{leaveEvent — odhlásenie z eventu:}
Pri odhlásení sa automaticky povýši prvý používateľ z poradovníka. Ak sa
používateľ odhlási po uplynutí deadlinu, je mu zaznamenaný no-show:

\begin{lstlisting}[language=Kotlin, caption=Odhlásenie z eventu]
fun leaveEvent(eventId: String, userId: String) {
    val applyPenalty = event.isPastCancelDeadline()
    var updated = event.copy(attendees = event.attendees - userId)
    if (updated.waitlist.isNotEmpty()) {
        val next = updated.waitlist.first()
        updated = updated.copy(
            attendees = updated.attendees + next,
            waitlist  = updated.waitlist.drop(1)
        )
    }
    if (applyPenalty) recordNoShow(userId, eventId)
}
\end{lstlisting}

\paragraph{removeAttendee — odstránenie účastníka tvorcom:}
Odstránený účastník je zároveň zablokovaný pre daný event a prvý z poradovníka
je automaticky povýšený:

\begin{lstlisting}[language=Kotlin, caption=Odstránenie účastníka tvorcom eventu]
fun removeAttendee(eventId: String, attendeeId: String) {
    val updated = event.copy(
        attendees   = event.attendees   - attendeeId,
        bannedUsers = event.bannedUsers + attendeeId
    )
    // + automatické povýšenie z poradovníka
}
\end{lstlisting}

% -------------------------------------------------------
\subsubsection{Filtrovanie udalostí}

Filtrovanie je implementované ako samostatná servisná trieda
\texttt{EventFilterService}, ktorá prijíma kritériá a aplikuje ich na zoznam
eventov.

\paragraph{FilterCriteria:}
Dátová trieda zapuzdruje všetky dostupné filtre:

\begin{lstlisting}[language=Kotlin, caption=Dátová trieda FilterCriteria]
data class FilterCriteria(
    val name           : String? = null,
    val category       : String? = null,
    val subcategory    : String? = null,
    val participants   : Int?    = null,
    val maxPrice       : Double? = null,
    val dateFrom       : Date?   = null,
    val dateTo         : Date?   = null,
    val maxDistanceKm  : Double? = null,
    val refLat         : Double  = 0.0,
    val refLng         : Double  = 0.0
)
\end{lstlisting}

\paragraph{EventFilterService:}
Trieda aplikuje kritériá postupne — ak event nevyhovuje čo i len jednému
z nich, je zo zoznamu vyradený:

\begin{lstlisting}[language=Kotlin, caption=Implementácia EventFilterService]
class EventFilterService(private val criteria: FilterCriteria) {

    fun filter(allEvents: List<Event>): List<Event> =
        allEvents.filter { matches(it) }

    private fun matches(event: Event): Boolean {
        if (!criteria.name.isNullOrEmpty() &&
            !event.title.lowercase().contains(criteria.name.lowercase()))
            return false

        if (!criteria.category.isNullOrEmpty() &&
            event.category != criteria.category)
            return false

        if (!criteria.subcategory.isNullOrEmpty() &&
            event.subcategory != criteria.subcategory)
            return false

        if (criteria.participants != null && criteria.participants > 0 &&
            event.participants < criteria.participants)
            return false

        if (criteria.maxPrice != null && criteria.maxPrice < 500.0 &&
            event.price > criteria.maxPrice)
            return false

        if (criteria.dateFrom != null && event.dateFrom != null &&
            event.dateFrom.before(criteria.dateFrom))
            return false

        if (criteria.maxDistanceKm != null) {
            val dist = haversineKm(
                criteria.refLat, criteria.refLng,
                event.latitude,  event.longitude
            )
            if (dist > criteria.maxDistanceKm) return false
        }
        return true
    }
}
\end{lstlisting}

\paragraph{Výpočet vzdialenosti — Haversine:}
Na výpočet vzdialenosti medzi dvoma GPS súradnicami sa používa Haversine vzorec,
ktorý zohľadňuje zakrivenie Zeme. Výsledkom je vzdialenosť v kilometroch
(polomer Zeme $R = 6371\ \text{km}$):

\begin{lstlisting}[language=Kotlin, caption=Haversine vzorec]
private fun haversineKm(
    lat1: Double, lon1: Double,
    lat2: Double, lon2: Double
): Double {
    val dLat = Math.toRadians(lat2 - lat1)
    val dLon = Math.toRadians(lon2 - lon1)
    val a = sin(dLat / 2).pow(2) +
            cos(Math.toRadians(lat1)) *
            cos(Math.toRadians(lat2)) *
            sin(dLon / 2).pow(2)
    return 6371.0 * 2 * asin(sqrt(a))
}
\end{lstlisting}

\paragraph{Filter bottom sheet — UI:}
Filter sa zobrazí ako \texttt{ModalBottomSheet} po kliknutí na tlačidlo
„Nájsť udalosť". Používateľ môže filtrovať podľa:
\begin{itemize}
    \item názvu — fulltextové vyhľadávanie (case-insensitive),
    \item kategórie a podkategórie — výber z \texttt{CATEGORY\_MAP},
    \item počtu osôb — slider v rozsahu 0–200,
    \item maximálnej ceny — slider 0–150\,€; hodnota nad 150\,€ znamená
          bez obmedzenia,
    \item vzdialenosti — slider 1–50\,km od GPS polohy používateľa;
          hodnota nad 50\,km znamená kdekoľvek,
    \item dátumu začiatku — výber cez \texttt{DatePickerDialog}
          a \texttt{TimePickerDialog}.
\end{itemize}

Po kliknutí na „Použiť" sa zostavia kritériá a filtrujú sa markery na mape:

\begin{lstlisting}[language=Kotlin, caption=Aplikovanie filtrov na mapu]
val criteria = FilterCriteria(
    name          = name.trim().ifEmpty { null },
    category      = selectedCategory.ifEmpty { null },
    maxPrice      = if (maxPrice >= 150f) Double.MAX_VALUE
                    else maxPrice.toDouble(),
    maxDistanceKm = if (maxDistance < 50f) maxDistance.toDouble()
                    else null,
    refLat        = userLocation?.latitude  ?: 0.0,
    refLng        = userLocation?.longitude ?: 0.0
)
val filtered = EventFilterService(criteria).filter(events)
markers.values.forEach { it.isVisible = false }
filtered.forEach { event -> markers[event.id]?.isVisible = true }
\end{lstlisting}

Výsledok filtra sa prejaví priamo na mape — zobrazia sa len markery eventov
spĺňajúcich všetky zadané kritériá. Stav filtra sa ukladá do
\texttt{AppViewModel.filterState}, čím sa zachováva pri navigácii medzi
obrazovkami.